\documentclass[11pt,a4paper]{article}
\usepackage[T1]{fontenc}
\usepackage{lmodern,amsmath,amssymb,amsthm,mathtools,booktabs,microtype}
\usepackage[margin=27mm]{geometry}
\usepackage[colorlinks=true,linkcolor=blue,citecolor=blue,urlcolor=blue]{hyperref}
\newtheorem{theorem}{Theorem}[section]
\newtheorem{proposition}[theorem]{Proposition}
\newtheorem{lemma}[theorem]{Lemma}
\newtheorem{corollary}[theorem]{Corollary}
\theoremstyle{remark}\newtheorem{remark}[theorem]{Remark}
\newcommand{\Q}{\mathbb Q}\newcommand{\Z}{\mathbb Z}
\newcommand{\R}{\mathbb R}
\DeclareMathOperator{\rad}{rad}
\DeclareMathOperator{\sfpart}{sf}
\setlength{\emergencystretch}{2em}
\title{Input Height Bounds and Odd Contact Radicals}
\author{}\date{10 November 2026}
\begin{document}\maketitle

\begin{abstract}
A positive power of a rational input gives a height scale whose
growth is uniform in an arbitrarily large additive integer shift.
We prove an exact inequality with factor $1/r$ for a power $r$,
and give a criterion that combines it with an elliptic index-height
estimate. Under rational $abc$, six fixed rational contacts and a
quadratic input scale give the required strict height inequality;
the precise absorption coefficient is $8/9$ for an explicit choice
of exponent. A cubic input scale also gives a height criterion for
five contacts. Four contacts behave differently: an explicit
elliptic curve supplies rational arguments of unbounded height
whose odd contact radical is one outside $\{2,3,5\}$. This last
statement is unconditional. All point arithmetic and displayed
polynomial identities in the example have an exact integer and
rational certificate.
\end{abstract}

\section{Heights and contact radicals}

For a reduced fraction $a/b$, with $b>0$, put
\[
                 h(a/b)=\log\max(|a|,b).
\]
For $u,v\in\Q$, write $u=A/D$, $v=C/D$ with
$D>0$ and $\gcd(D,A,C)=1$, and define
\[
                 h(1:u:v)=\log\max(D,|A|,|C|).
\]
Both $h(u)$ and $h(v)$ are at most this projective height.
For a nonzero integer $m$, let $\rad(m)$ be the product of
its distinct prime divisors, and let $\sfpart(m)$ be the
product of primes occurring to odd exponent. Signs are omitted.
For a prime $p$, the valuation $v_p$ is normalized by $v_p(p)=1$.

Let $\Theta\subset\Q$ be a finite set of distinct contacts.
Choose a finite set $S$ of primes containing the denominators of
the contacts and all primes at which two contacts have the same
reduction. For $t\notin\Theta$, put
\begin{equation}\label{eq:radical}
 M_{\Theta,S}(t)=
 \prod_{\substack{p\notin S:\text{ some }\theta\in\Theta\text{ has}\\
                 v_p(t-\theta)>0\text{ odd}}}p.
\end{equation}
The empty product is one. Enlarging $S$ changes a logarithmic
bound for this radical by at most a fixed additive constant.

We use the following form of rational $abc$: for every
$\epsilon>0$ there is $K_\epsilon\ge1$ such that
\[
 c\le K_\epsilon\rad(abc)^{1+\epsilon}
 \quad(a+b=c, a,b,c>0, \gcd(a,b)=1).
\]
When a single exponent $\kappa>1$ is sufficient, we state that
we assume a constant $K_\kappa\ge1$ with
\begin{equation}\label{eq:weakabc}
 c\le K_\kappa\rad(abc)^\kappa
 \quad(a+b=c, a,b,c>0, \gcd(a,b)=1).
\end{equation}

\section{A height scale uniform in its additive shift}

\begin{lemma}\label{lem:power}
Let $r\ge1$ and $B\ge2$ be integers. Suppose that
$\gamma\in\Q$ satisfies either $\gamma\ge1$ or that $r$ is even.
Set
\[
                         \beta=2\gamma^r+B.
\]
Then $\beta\ge B$ and
\begin{equation}\label{eq:powerheight}
                         h(\gamma)\le\frac1r h(\beta).
\end{equation}
In particular the inequality is uniform in $B$.
\end{lemma}
\begin{proof}
Write $\gamma=a/b$ with $\gcd(a,b)=1$ and $b>0$.
The cancellation factor is
\[
 g=\gcd(2a^r+Bb^r,b^r)=\gcd(2,b^r)\le2.
\]
Under either stated condition, the numerator is positive and
\[
                   2a^r+Bb^r\ge2\max(|a|^r,b^r).
\]
It is also at least the denominator. Consequently
\[
 h(\beta)=\log\frac{2a^r+Bb^r}{g}
                \ge r\log\max(|a|,b)=r h(\gamma).
\]
\end{proof}

Two input choices will be useful. The direct choice
$\gamma=X$ permits every rational $X$ when $r$ is even.
The positive choice
\begin{equation}\label{eq:positiveinput}
                         \gamma=2X^2+1
\end{equation}
permits every $r\ge1$. For rational $X$, it satisfies
\begin{equation}\label{eq:integrality}
                      \gamma\in\Z\quad\Longleftrightarrow\quad X\in\Z.
\end{equation}
Indeed, if $X=c/d$ is reduced, integrality of $2X^2+1$ gives
$d^2\mid2$, hence $d=1$. For integer input, $\gamma$ is a
positive odd integer. If $B$ is odd, then $2\gamma^r+B$ is
also a positive odd integer.

\begin{theorem}[A uniform bound from a smaller input height]
\label{thm:absorption}
Let $T:\Z\setminus\{0\}\to\Q$ satisfy
\begin{equation}\label{eq:indexheight}
                  n^2\le A h(T(n))+J
\end{equation}
for fixed $A>0$ and $J\ge0$. Let $r\ge1$, let $B\ge2$ be
an integer, and put $\beta=2\gamma^r+B$ under the hypotheses
of Lemma~\ref{lem:power}. Set
\[
                      H=h(1:\beta:T(n)).
\]
Suppose that, whenever $\gamma\ne n$, the arithmetic conditions
under consideration imply
\begin{equation}\label{eq:residualheight}
               H\le\kappa\log N_\gamma+C,
               \qquad 0<\kappa<r,
\end{equation}
where $N_\gamma$ is the positive absolute reduced numerator of
$\gamma-n$ and $C\ge0$ is independent of $B,\gamma,n$.
Then $|n|$ and $H$ are bounded independently of $B$ whenever
$\gamma\ne n$. In particular, for every sufficiently large $B$,
these conditions imply $\gamma=n$.
\end{theorem}
\begin{proof}
If $\gamma=a/b$ is reduced with $b>0$, then
$\gcd(a-nb,b)=1$ and
\[
 N_\gamma=|a-nb|\le e^{h(\gamma)}(1+|n|).
\]
By Lemma~\ref{lem:power},
\[
                  \log N_\gamma\le H/r+\log(1+|n|).
\]
Insert this into~\eqref{eq:residualheight} and move the term
$(\kappa/r)H$ to the left. With
\[
 c=\frac{\kappa}{1-\kappa/r},\qquad
 C'=\frac{C}{1-\kappa/r},
\]
one obtains
\begin{equation}\label{eq:absorbed}
                         H\le c\log(1+|n|)+C'.
\end{equation}
Since $h(T(n))\le H$, equation~\eqref{eq:indexheight} gives
\[
 n^2\le Ac\log(1+|n|)+AC'+J\le Ac|n|+AC'+J.
\]
Thus any integer $M>Ac+\sqrt{AC'+J}$ bounds $|n|$ strictly.
It follows that
\[
 H\le c\log(1+M)+C'=:H_*.
\]
But $H\ge h(\beta)\ge\log\beta\ge\log B$.
Choosing $B>e^{H_*}$ rules out $\gamma\ne n$.
\end{proof}

\begin{remark}
For the positive input~\eqref{eq:positiveinput}, the conclusion
$\gamma=n$ proves $X\in\Z$ by~\eqref{eq:integrality}; the
completeness label $n=2X^2+1$ never vanishes. For the direct
input $\gamma=X$, the conclusion is already $X\in\Z$.
If a polynomial $Q(X,\boldsymbol Y)$ describes precisely the
nonzero integers, then $XQ(X,\boldsymbol Y)$ describes all
integers, with the same variables and one higher degree.
This last assertion is the elementary equivalence
$XQ=0\Longleftrightarrow X=0\text{ or }Q=0$ over $\Q$.
\end{remark}

\section{Contact and power-sum height estimates}

\begin{proposition}\label{prop:contacts}
Assume rational $abc$. Let $\Theta$ consist of $d\ge5$ distinct
rational contacts, and let $S$ be as in~\eqref{eq:radical}.
For every $\epsilon$ with $0<2\epsilon<d-4$, there is a
constant $C_{\Theta,S,\epsilon}\ge0$ such that
\begin{equation}\label{eq:contactheight}
 h(t)\le\frac{1}{d-4-2\epsilon}\log M_{\Theta,S}(t)
                   +C_{\Theta,S,\epsilon}
 \qquad(t\in\Q\setminus\Theta).
\end{equation}
\end{proposition}
\begin{proof}
Write each contact as $a_i/b_i$ in lowest terms with $b_i>0$,
and set
\[
                   G(U,V)=\prod_{i=1}^d(b_iU-a_iV).
\]
This integral homogeneous form has no repeated factor.
For a primitive integer pair $(u,v)$ with $v>0$, put
$H_0=\max(|u|,v)$. Granville's consequence of rational
$abc$~\cite[Proposition~1]{Granville} gives
\[
             \log\rad(G(u,v))\ge(d-2-\epsilon)\log H_0-C_1.
\]
The finitely many bounded heights can be included by enlarging
$C_1$. Also $|G(u,v)|\le C_2H_0^d$. The primewise inequality
\[
                       \rad(m)^2\le |m|\sfpart(m)
\]
therefore gives
\[
 \log\sfpart(G(u,v))\ge(d-4-2\epsilon)\log H_0-C_3.
\]
Outside $S$, at most one primitive linear factor can be
divisible by a given prime, and a prime dividing $v$ divides
none. Its odd exponent occurs exactly at a positive odd
valuation of one of $u/v-a_i/b_i$. Hence
\[
 \sfpart(G(u,v))\le\left(\prod_{p\in S}p\right)M_{\Theta,S}(u/v).
\]
Dividing the resulting inequality by $d-4-2\epsilon>0$
proves~\eqref{eq:contactheight}.
\end{proof}

\begin{lemma}\label{lem:fermat}
Assume~\eqref{eq:weakabc} for a fixed $\kappa>1$.
Let $m\ge4$ be even and put $\alpha=m/\kappa-3$.
If $\alpha>0$, then
\begin{equation}\label{eq:fermatheight}
 h(1:u:v)\le\alpha^{-1}h(u^m+v^m)+C_{m,\kappa}
 \qquad(u,v\in\Q),
\end{equation}
where one may take
\[
 C_{m,\kappa}=\frac{\alpha+1}{m\alpha}\log2
                       +\frac{1}{\kappa\alpha}\log K_\kappa.
\]
\end{lemma}
\begin{proof}
If a coordinate is zero, the left side is $h(u^m+v^m)/m$;
if both are zero it is zero. These cases satisfy the assertion.
Otherwise write
\[
 u=gA/D,\qquad v=gC/D,
 \qquad \gcd(A,C)=\gcd(g,D)=1,\quad g,D>0.
\]
Put $H_0=\max(|A|,|C|)$ and $N=A^m+C^m$, so
$H_0^m\le N\le2H_0^m$. Set
\[
 k=\gcd(N,D^m),\qquad m_0=N/k,\qquad n_0=D^m/k,
 \qquad M=\max(g^m m_0,n_0).
\]
Then $u^m+v^m=g^m m_0/n_0$ is reduced. Every prime dividing
$N$ either divides $D$ or survives in $m_0$, and so
$\rad(N)\le Dm_0$. Applying~\eqref{eq:weakabc} to the
pairwise coprime positive integers $A^m,C^m,N$ gives
\[
 H_0^{m/\kappa-2}\le K_\kappa^{1/\kappa}Dm_0.
\]
Since $D^m=n_0N/m_0\le2n_0H_0^m/m_0$, this implies
\[
 H_0^\alpha\le2^{1/m}K_\kappa^{1/\kappa}
                         m_0^{(m-1)/m}n_0^{1/m}.
\]
Moreover
\[
 \max(gH_0,D)\le2^{1/m}H_0(M/m_0)^{1/m}.
\]
Combining the logarithms of these inequalities gives
\[
 h(1:u:v)\le C_{m,\kappa}
 +\frac{m-1-\alpha}{m\alpha}\log m_0
 +\frac{1}{m\alpha}\log n_0+\frac1m\log M.
\]
All three coefficients are nonnegative and sum to $1/\alpha$.
Since $m_0,n_0\le M$, this proves~\eqref{eq:fermatheight}.
\end{proof}

Rational $abc$ supplies $\kappa=16/15$. With $m=4$, the
lemma gives
\begin{equation}\label{eq:fourthheight}
 h(1:u:v)\le\frac43h(u^4+v^4)+C_4,
 \qquad C_4=\frac7{12}\log2+\frac54\log K_{16/15}.
\end{equation}

\section{Six contacts and a quadratic input scale}

\begin{theorem}\label{thm:six}
Assume rational $abc$ and the index-height bound
\eqref{eq:indexheight}. Fix six distinct rational contacts and
the associated radical $M=M_{\Theta,S}$. Put
\[
 \gamma=2X^2+1,\qquad \beta=2\gamma^2+B,
 \qquad a_1=\beta^4+T(n)^4,
\]
where $B\ge3$ is an odd fixed integer and $n\ne0$ is an
integer label. Let $\lambda\in\Q^\times$. Suppose that, when $\gamma\ne n$, the local
conditions give
\begin{equation}\label{eq:twodivisibilities}
             M(\lambda)\mid N_\gamma,
             \qquad M(a_1\lambda)\mid N_\gamma,
\end{equation}
with both arguments outside the contact set. Then every
sufficiently large odd $B$ forces $X\in\Z$.
The sufficient bound is uniform in the input and witnesses.
\end{theorem}
\begin{proof}
In Proposition~\ref{prop:contacts}, take $d=6$ and
$\epsilon=1/4$. Then
\[
                       h(t)\le\frac23\log M(t)+C_0.
\]
By~\eqref{eq:twodivisibilities}, the quotient-height inequality
and~\eqref{eq:fourthheight} give, with $H=h(1:\beta:T(n))$,
\[
 \begin{split}
 H&\le\frac43h(a_1)+C_4\\
  &\le\frac43\bigl(h(\lambda)+h(a_1\lambda)\bigr)+C_4\\
  &\le\frac{16}{9}\log N_\gamma+C_2,
       \qquad C_2=\frac83C_0+C_4.
 \end{split}
\]
Lemma~\ref{lem:power} yields $h(\gamma)\le H/2$. Consequently
\[
 H\le\frac89H+\frac{16}{9}\log(1+|n|)+C_2,
 \qquad H\le16\log(1+|n|)+9C_2.
\]
Theorem~\ref{thm:absorption}, with $r=2$ and $\kappa=16/9$,
now gives $\gamma=n$ for sufficiently large $B$.
Equation~\eqref{eq:integrality} proves the assertion.
\end{proof}

The same calculation gives other height criteria. The table
records choices in Proposition~\ref{prop:contacts}, using the
fourth-power estimate~\eqref{eq:fourthheight}. The column
headed $\kappa/r$ is the coefficient of $H$ that is moved to
the left in the proof of Theorem~\ref{thm:absorption}.
\begin{center}
\begin{tabular}{ccccc}
\toprule
Contacts $d$ & $\epsilon$ & Contact coefficient & Input power $r$ & $\kappa/r$\\
\midrule
8 & $1/2$ & $1/3$ & 1 & $8/9$\\
6 & $1/4$ & $2/3$ & 2 & $8/9$\\
5 & $1/36$ & $18/17$ & 3 & $16/17$\\
\bottomrule
\end{tabular}
\end{center}
For example, five contacts and $\beta=2\gamma^3+B$ give
\[
 H\le\frac{48}{17}\log N_\gamma+C
       \le\frac{16}{17}H+\frac{48}{17}\log(1+|n|)+C,
\]
and hence $H\le48\log(1+|n|)+17C$.
These assertions are height implications under the displayed
local divisibilities. A polynomial construction must separately
supply those divisibilities and prove completeness of its local
and pole conditions.

\begin{remark}
The smaller input height helps only when the comparison residual
is $\gamma-n$. Keeping the residual $\beta-n$ gives
$\log N_\beta\le h(\beta)+\log(1+|n|)$, which does not acquire
the factor $1/r$ from Lemma~\ref{lem:power}.
\end{remark}

\section{Four contacts with bounded odd radical}

The degree restriction in Proposition~\ref{prop:contacts} reflects
a genuine obstruction. For four contacts its radical argument
has exponent $-2\epsilon$ after the square part is removed.
The following explicit example shows that no replacement
logarithmic estimate can hold for this contact set.

\begin{theorem}\label{thm:fourcounterexample}
Let
\[
                  \Theta=\{-1,0,2,3\},\qquad S\supseteq\{2,3,5\}
\]
with $S$ finite. There are infinitely many
$t\in\Q\setminus(\Theta\cup\{1\})$ for which
\[
                         M_{\Theta,S}(t)=1.
\]
Their heights are unbounded. Consequently there are no constants
$c,C\ge0$ such that
\[
                h(t)\le c\log M_{\Theta,S}(t)+C
\]
for every rational $t$ outside the contacts and the pole.
No conjecture is used.
\end{theorem}
\begin{proof}
Consider the smooth quartic curve
\begin{equation}\label{eq:quartictwist}
                   10y^2=(t+1)t(t-2)(t-3).
\end{equation}
Its four roots are distinct, and it has the rational point
$(t,y)=(4,2)$. On $t\ne0$, make the substitutions
\[
                        X_E=\frac{60}{t},\qquad
                        Y_E=\frac{600y}{t^2}.
\]
They give the elliptic curve
\begin{equation}\label{eq:longcurve}
 E:\quad Y_E^2=X_E^3+10X_E^2-2400X_E+36000
              =(X_E+60)(X_E-20)(X_E-30),
\end{equation}
with inverse formulas
\[
                         t=\frac{60}{X_E},\qquad
                         y=\frac{6Y_E}{X_E^2}.
\]
The point $(4,2)$ corresponds to $P=(15,75)$.
The tangent slope at $P$ is $-19/2$, and the group law gives
\begin{equation}\label{eq:double}
                         2P=\left(\frac{201}{4},\frac{2079}{8}\right).
\end{equation}

To apply the integral short-form version of Nagell--Lutz, put
\[
                        x=9X_E+30,\qquad Y=27Y_E.
\]
Then
\begin{equation}\label{eq:shortcurve}
                 Y^2=x^3-197100x+32130000.
\end{equation}
This integral equation is nonsingular; its cubic has the distinct
roots $-510,210,300$. The image of $2P$ has abscissa $1929/4$,
which is not integral. Nagell--Lutz states that every rational
torsion point on an integral short Weierstrass equation has
integral coordinates; see~\cite[Theorem~1.1]{Alpoge} or
\cite[Chapter~II, Theorem~5.1]{Milne}. Thus $P$ has infinite order.

For each nonzero integer $n$, the point $[n]P$ is nonidentity
and is not a point of order two. It cannot have $X_E=0$, since
$Y_E^2=36000$ has no rational solution. The inverse formulas
therefore give a rational point $(t_n,y_n)$ on
\eqref{eq:quartictwist} with $t_n\notin\Theta$.
Also $t_n\ne1$, since at $t=1$ the equation would require
$y^2=2/5$. There are infinitely many distinct $t_n$: a fixed
value of $t$ determines at most two points of~\eqref{eq:longcurve},
whereas the multiples of $P$ are all distinct.

Fix $p\notin\{2,3,5\}$. If $v_p(t_n)<0$, then all four
differences $t_n-\theta$ have negative valuation, so $p$
does not occur in~\eqref{eq:radical}. If $v_p(t_n)\ge0$,
at most one difference has positive valuation, since the
nonzero differences among the contacts have prime factors
only two and three. Equation~\eqref{eq:quartictwist} gives
\[
 v_p\bigl((t_n+1)t_n(t_n-2)(t_n-3)\bigr)=2v_p(y_n),
\]
which is even. Therefore no positive contact valuation at such
a prime is odd. It follows that $M_{\Theta,S}(t_n)=1$.

Finally, only finitely many reduced rational fractions have
height at most any fixed constant. Since the $t_n$ form an
infinite set, their heights are unbounded. This proves all
assertions.
\end{proof}

\begin{corollary}\label{cor:positive}
The counterexamples in Theorem~\ref{thm:fourcounterexample}
may be chosen with $3<t<5$, and still have unbounded height.
\end{corollary}
\begin{proof}
The real group $E(\R)$ is compact. The closure of $\Z P$
is an infinite compact subgroup, hence has no isolated point:
otherwise it would be discrete and compact, and therefore finite.
Every sufficiently small neighborhood of $P=(15,75)$ consequently
contains infinitely many multiples of $P$. The rational function
$60/X_E$ is continuous there and equals four at $P$.
Choose the neighborhood so that its values lie in $(3,5)$.
The preceding finite-fiber and height arguments apply to this
infinite subset.
\end{proof}

The theorem concerns a proposed height bound for the odd radical.
It does not assert that these arguments extend to complete
nonintegral witnesses of any Diophantine construction.

\section{Exact certificate}

The standard-library program
\nolinkurl{computation/contact_heights/check_input_height_amplification.py}
checks the rational point and duplication in~\eqref{eq:longcurve},
the birational substitution identity, the integral change to
\eqref{eq:shortcurve}, and the nonintegral doubled abscissa.
It also verifies the distinct roots, the excluded rational
values, and the exact rational constants in the height table.
Its machine-readable result is
\nolinkurl{input-height-amplification-certificate.json} in the
same directory. The infinitude assertion follows from the
proved nontorsion criterion and the argument above; it is not
inferred from a finite search over multiples.

\begin{thebibliography}{9}
\bibitem{Granville}
A.~Granville, Rational and integral points on quadratic twists of a
given hyperelliptic curve, \emph{International Mathematics Research
Notices} (2007), article rnm027, Proposition~1,
\url{https://doi.org/10.1093/imrn/rnm027}; author manuscript,
\url{https://dms.umontreal.ca/~andrew/PDF/hyperelliptics.pdf}.
\bibitem{Alpoge}
L.~Alpoge, \emph{Nagell--Lutz, quickly}, author manuscript,
Theorem~1.1,
\url{https://people.math.harvard.edu/~alpoge/papers/nagell-lutz,%20quickly.pdf}.
\bibitem{Milne}
J.~S.~Milne, \emph{Elliptic Curves}, second edition,
World Scientific, 2021, Chapter~II, Section~5,
\url{https://www.jmilne.org/math/Books/EC2.pdf}.
\end{thebibliography}
\end{document}
